\section{Wavefunction Description}
\label{sec:wavefunction}

This section develops the quantum mechanical description of I-particles through wavefunction formalism.

\subsection{Quantum Mechanical Wavefunction}
\label{subsec:quantum_wavefunction}

In quantum mechanical treatments, a single I particle in a gravitational or potential well can be described by a wavefunction:

\begin{equation}
\label{eq:i_wavefunction}
\Psi_I(\mathbf{r}, t) = A e^{i(\mathbf{k} \cdot \mathbf{r} - \omega t)}
\end{equation}

where:
\begin{itemize}
    \item $A$ is the amplitude (normalized by boundary or normalization conditions)
    \item $\mathbf{k}$ is the wavevector, proportional to the momentum of the particle: $\mathbf{p} = \hbar\mathbf{k}$
    \item $\omega$ is the angular frequency, related to the particle's energy by $E = \hbar\omega$
\end{itemize}

This form implies that I exhibits wave-particle duality, typical of quantum entities.

\subsubsection{Normalization Condition}

The wavefunction must satisfy:

\begin{equation}
\label{eq:normalization}
\int_{\text{all space}} |\Psi_I(\mathbf{r}, t)|^2 \, d^3r = 1
\end{equation}

For the plane wave \eqref{eq:i_wavefunction}, we use box normalization with volume $V$:

\begin{equation}
\label{eq:box_normalization}
A = \frac{1}{\sqrt{V}}
\end{equation}

\subsubsection{Dispersion Relations}

For free I-particles:

\begin{align}
\label{eq:dispersion_relativistic}
\omega &= c\sqrt{k^2 + (m_I c/\hbar)^2} \quad \text{(relativistic)} \\
\label{eq:dispersion_nonrelativistic}
\omega &= \frac{\hbar k^2}{2m_I} \quad \text{(non-relativistic, } k \ll m_I c/\hbar\text{)}
\end{align}

\subsection{Schrödinger Equation}

The time evolution of the I-wavefunction is governed by the Schrödinger equation:

\begin{equation}
\label{eq:schrodinger}
i\hbar\frac{\partial\Psi_I}{\partial t} = \hat{H}_I \Psi_I
\end{equation}

where $\hat{H}_I$ is the Hamiltonian operator.

\subsubsection{Free Particle Hamiltonian}

For a free I-particle:

\begin{equation}
\label{eq:free_hamiltonian}
\hat{H}_I = -\frac{\hbar^2}{2m_I}\nabla^2 + m_I c^2
\end{equation}

\subsubsection{Particle in a Potential}

In the presence of a potential $V(\mathbf{r})$:

\begin{equation}
\label{eq:hamiltonian_potential}
\hat{H}_I = -\frac{\hbar^2}{2m_I}\nabla^2 + V(\mathbf{r})
\end{equation}

Time-independent solutions satisfy:

\begin{equation}
\label{eq:time_independent_schrodinger}
\hat{H}_I \psi_n(\mathbf{r}) = E_n \psi_n(\mathbf{r})
\end{equation}

\subsection{Wavefunction in Gravitational Wells}

In a gravitational or potential well, $\omega$ and $\mathbf{k}$ may vary with position if the potential is non-uniform.

\subsubsection{WKB Approximation}

For slowly varying potentials, the WKB (Wentzel-Kramers-Brillouin) approximation:

\begin{equation}
\label{eq:wkb_wavefunction}
\Psi_I(\mathbf{r}) \approx \frac{A}{\sqrt{k(\mathbf{r})}} \exp\left(i\int_{\mathbf{r}_0}^{\mathbf{r}} \mathbf{k}(\mathbf{r}') \cdot d\mathbf{r}'\right)
\end{equation}

where the local wavevector:

\begin{equation}
\label{eq:local_wavevector}
k(\mathbf{r}) = \frac{1}{\hbar}\sqrt{2m_I[E - V(\mathbf{r})]}
\end{equation}

\subsubsection{Gravitational Potential Well}

For I-particles in the gravitational potential of a massive object:

\begin{equation}
\label{eq:gravitational_well}
V(r) = -\frac{Gm_I M}{r}
\end{equation}

The radial wavefunction satisfies:

\begin{equation}
\label{eq:radial_schrodinger}
-\frac{\hbar^2}{2m_I}\frac{d^2u}{dr^2} + \left[V(r) + \frac{\hbar^2\ell(\ell+1)}{2m_I r^2}\right]u = E u
\end{equation}

where $u(r) = r R(r)$ and $\ell$ is the angular momentum quantum number.

\subsubsection{Bound State Solutions}

For bound states ($E < 0$), the energy levels are quantized:

\begin{equation}
\label{eq:bound_state_energy}
E_n = -\frac{G^2m_I^3M^2}{2\hbar^2 n^2} \quad (n = 1, 2, 3, \ldots)
\end{equation}

This is analogous to the hydrogen atom with:

\begin{equation}
\label{eq:gravitational_bohr_radius}
a_G = \frac{\hbar^2}{Gm_I^2 M}
\end{equation}

\subsection{Many-Body Wavefunction}

For a system of $N$ I-particles:

\begin{equation}
\label{eq:many_body_wavefunction}
\Psi_I(\mathbf{r}_1, \mathbf{r}_2, \ldots, \mathbf{r}_N, t)
\end{equation}

\subsubsection{Symmetrization}

If I-particles are bosons (integer spin):

\begin{equation}
\label{eq:symmetric_wavefunction}
\Psi_I(\ldots, \mathbf{r}_i, \ldots, \mathbf{r}_j, \ldots) = +\Psi_I(\ldots, \mathbf{r}_j, \ldots, \mathbf{r}_i, \ldots)
\end{equation}

If I-particles are fermions (half-integer spin):

\begin{equation}
\label{eq:antisymmetric_wavefunction}
\Psi_I(\ldots, \mathbf{r}_i, \ldots, \mathbf{r}_j, \ldots) = -\Psi_I(\ldots, \mathbf{r}_j, \ldots, \mathbf{r}_i, \ldots)
\end{equation}

\subsubsection{Second Quantization}

The field operator expansion:

\begin{equation}
\label{eq:field_operator}
\hat{\Psi}_I(\mathbf{r}, t) = \sum_\mathbf{k} \frac{1}{\sqrt{V}} e^{i\mathbf{k}\cdot\mathbf{r}} \hat{a}_\mathbf{k}(t)
\end{equation}

For bosons:
\begin{equation}
\label{eq:boson_commutation}
[\hat{a}_\mathbf{k}, \hat{a}_{\mathbf{k}'}^\dagger] = \delta_{\mathbf{k}\mathbf{k}'}
\end{equation}

For fermions:
\begin{equation}
\label{eq:fermion_anticommutation}
\{\hat{a}_\mathbf{k}, \hat{a}_{\mathbf{k}'}^\dagger\} = \delta_{\mathbf{k}\mathbf{k}'}
\end{equation}

\subsection{Probability Current and Conservation}

The probability density:

\begin{equation}
\label{eq:probability_density}
\rho(\mathbf{r}, t) = |\Psi_I(\mathbf{r}, t)|^2 = \Psi_I^* \Psi_I
\end{equation}

The probability current density:

\begin{equation}
\label{eq:probability_current}
\mathbf{j}(\mathbf{r}, t) = \frac{\hbar}{2m_I i}(\Psi_I^*\nabla\Psi_I - \Psi_I\nabla\Psi_I^*)
\end{equation}

Continuity equation:

\begin{equation}
\label{eq:continuity}
\frac{\partial\rho}{\partial t} + \nabla \cdot \mathbf{j} = 0
\end{equation}

\subsubsection{Current for Plane Waves}

For $\Psi_I = A e^{i(\mathbf{k}\cdot\mathbf{r} - \omega t)}$:

\begin{equation}
\label{eq:plane_wave_current}
\mathbf{j} = |A|^2 \frac{\hbar\mathbf{k}}{m_I} = \rho \mathbf{v}
\end{equation}

where $\mathbf{v} = \hbar\mathbf{k}/m_I$ is the particle velocity.

\subsection{Coherent States}

Coherent states of I-fields:

\begin{equation}
\label{eq:coherent_state}
|\alpha\rangle_I = e^{-|\alpha|^2/2} \sum_{n=0}^\infty \frac{\alpha^n}{\sqrt{n!}} |n\rangle
\end{equation}

are eigenstates of the annihilation operator:

\begin{equation}
\label{eq:coherent_eigenvalue}
\hat{a}|\alpha\rangle = \alpha|\alpha\rangle
\end{equation}

\subsubsection{Coherent State Wavefunction}

In position space:

\begin{equation}
\label{eq:coherent_position}
\Psi_\alpha(\mathbf{r}) = \langle\mathbf{r}|\alpha\rangle = \left(\frac{m_I\omega}{\pi\hbar}\right)^{3/4} \exp\left(-\frac{m_I\omega}{2\hbar}(\mathbf{r}-\mathbf{r}_0)^2 + i\mathbf{p}_0\cdot\mathbf{r}/\hbar\right)
\end{equation}

\subsection{Entangled States}

Two-particle entangled states:

\begin{equation}
\label{eq:entangled_state}
|\Psi\rangle_{12} = \frac{1}{\sqrt{2}}(|0\rangle_1|1\rangle_2 + e^{i\phi}|1\rangle_1|0\rangle_2)
\end{equation}

The reduced density matrix:

\begin{equation}
\label{eq:reduced_density}
\rho_1 = \text{Tr}_2[|\Psi\rangle_{12}\langle\Psi|] = \frac{1}{2}(|0\rangle\langle 0| + |1\rangle\langle 1|)
\end{equation}

Entanglement entropy:

\begin{equation}
\label{eq:entanglement_entropy}
S_{\text{ent}} = -\text{Tr}(\rho_1 \ln \rho_1) = \ln 2
\end{equation}

\subsection{Time Evolution and Propagators}

\subsubsection{Free Particle Propagator}

The propagator for free I-particles:

\begin{equation}
\label{eq:free_propagator}
K(\mathbf{r}, t; \mathbf{r}_0, 0) = \left(\frac{m_I}{2\pi i\hbar t}\right)^{3/2} \exp\left(\frac{im_I(\mathbf{r}-\mathbf{r}_0)^2}{2\hbar t}\right)
\end{equation}

satisfies:

\begin{equation}
\label{eq:propagator_equation}
i\hbar\frac{\partial K}{\partial t} = -\frac{\hbar^2}{2m_I}\nabla^2 K
\end{equation}

\subsubsection{Path Integral Formulation}

The wavefunction can be expressed as:

\begin{equation}
\label{eq:path_integral}
\Psi_I(\mathbf{r}, t) = \int d^3r_0 \, K(\mathbf{r}, t; \mathbf{r}_0, 0) \Psi_I(\mathbf{r}_0, 0)
\end{equation}

where:

\begin{equation}
\label{eq:feynman_path_integral}
K(\mathbf{r}_f, t_f; \mathbf{r}_i, t_i) = \int_{\mathbf{r}_i}^{\mathbf{r}_f} \mathcal{D}[\mathbf{r}(t)] \, e^{iS[\mathbf{r}(t)]/\hbar}
\end{equation}

and $S[\mathbf{r}(t)]$ is the classical action.

\subsection{Wavefunction Collapse and Measurement}

Upon measurement of position:

\begin{equation}
\label{eq:position_measurement}
\Psi_I(\mathbf{r}, t) \xrightarrow{\text{measurement}} \delta(\mathbf{r} - \mathbf{r}_{\text{measured}})
\end{equation}

The probability of measuring position in volume $d^3r$:

\begin{equation}
\label{eq:measurement_probability}
dP = |\Psi_I(\mathbf{r}, t)|^2 d^3r
\end{equation}

\subsubsection{I-Decoherence Model}

In I-theory, measurement is described by I-decoherence:

\begin{equation}
\label{eq:i_decoherence_wavefunction}
\Psi_I(\mathbf{r}, t) = \sum_n c_n(t) \psi_n(\mathbf{r}) \xrightarrow{t\to\infty} \psi_m(\mathbf{r}) \oplus I\sum_{n\neq m} c_n\psi_n(\mathbf{r})
\end{equation}

The decoherence functional:

\begin{equation}
\label{eq:decoherence_functional}
\mathcal{D}[t] = \exp\left(-\int_0^t \Gamma_I(t') dt'\right)
\end{equation}

where $\Gamma_I$ is the I-decoherence rate.

\subsection{Wavefunction in Curved Spacetime}

In general relativistic contexts:

\begin{equation}
\label{eq:curved_spacetime_wavefunction}
\left(g^{\mu\nu}\nabla_\mu\nabla_\nu + \frac{m_I^2c^2}{\hbar^2}\right)\Psi_I = 0
\end{equation}

where $\nabla_\mu$ is the covariant derivative.

\subsubsection{Schwarzschild Spacetime}

Near a black hole with metric:

\begin{equation}
\label{eq:schwarzschild_metric}
ds^2 = -\left(1-\frac{2GM}{c^2r}\right)c^2dt^2 + \left(1-\frac{2GM}{c^2r}\right)^{-1}dr^2 + r^2d\Omega^2
\end{equation}

The effective potential for radial I-waves:

\begin{equation}
\label{eq:effective_potential_schwarzschild}
V_{\text{eff}}(r) = \left(1-\frac{2GM}{c^2r}\right)\left(\frac{\ell(\ell+1)\hbar^2}{2m_Ir^2} + \frac{m_Ic^2}{2}\right)
\end{equation}

\subsection{Applications}

\subsubsection{I-Particle Interference}

Double-slit experiment with I-particles:

\begin{equation}
\label{eq:double_slit}
\Psi_I = \Psi_1 + \Psi_2 = A(e^{i\mathbf{k}\cdot\mathbf{r}_1} + e^{i\mathbf{k}\cdot\mathbf{r}_2})
\end{equation}

Intensity pattern:

\begin{equation}
\label{eq:interference_pattern}
I(\mathbf{r}) = |\Psi_I|^2 = 2|A|^2[1 + \cos(\mathbf{k}\cdot\Delta\mathbf{r})]
\end{equation}

\subsubsection{Quantum Tunneling}

Tunneling probability through a barrier:

\begin{equation}
\label{eq:tunneling_probability}
T \approx \exp\left(-\frac{2}{\hbar}\int_{x_1}^{x_2}\sqrt{2m_I(V(x)-E)}\,dx\right)
\end{equation}

where $x_1$ and $x_2$ are the classical turning points.
